\documentclass[12pt,letterpaper]{article}
\usepackage[utf8]{inputenc}
\usepackage[T1]{fontenc}
\usepackage[spanish]{babel}
\usepackage{geometry}
\usepackage{enumitem}
\usepackage{array}
\usepackage{booktabs}
\usepackage{titlesec}
\usepackage{fancyhdr}
\usepackage{xcolor}
\usepackage{helvet}
\renewcommand{\familydefault}{\sfdefault}

\geometry{margin=2.5cm}

\definecolor{azulgob}{RGB}{31,56,100}
\definecolor{grissuave}{RGB}{100,100,100}

\titleformat{\section}{\large\bfseries\color{azulgob}}{}{0em}{}[\titlerule]
\titleformat{\subsection}{\normalsize\bfseries\color{azulgob}}{}{0em}{}
\titlespacing*{\section}{0pt}{20pt}{10pt}
\titlespacing*{\subsection}{0pt}{14pt}{6pt}

\pagestyle{fancy}
\fancyhf{}
\renewcommand{\headrulewidth}{0.4pt}
\fancyhead[L]{\small\color{grissuave}Casa de la Cultura del Municipio}
\fancyhead[R]{\small\color{grissuave}Preguntas para Reunión}
\fancyfoot[C]{\small\color{grissuave}\thepage}

\setlist[itemize]{leftmargin=1.5em, itemsep=8pt, parsep=2pt}

\begin{document}

% --- Portada ---
\begin{titlepage}
\centering
\vspace*{3cm}

{\color{azulgob}\rule{\textwidth}{2pt}}

\vspace{1.5cm}

{\Huge\bfseries\color{azulgob} Preguntas para Reunión}

\vspace{0.8cm}

{\Large Sistema de Registro y Capacitación\\de Servidores Públicos}

\vspace{1.5cm}

{\color{azulgob}\rule{\textwidth}{2pt}}

\vspace{2cm}

{\large Casa de la Cultura del Municipio}

\vspace{0.5cm}

{\large\color{grissuave} Área de Desarrollo de Sistemas}

\vspace{2cm}

{\large 21 de junio de 2026}

\vfill

{\small\color{grissuave} Documento elaborado para facilitar la toma de decisiones.\\Las respuestas permitirán alimentar el sistema y definir reglas operativas.}

\end{titlepage}

% --- Contenido ---

\section{I. Preguntas Prioritarias --- Alimentación de Datos}

\subsection{1. Catálogo de cursos de capacitación}

\begin{itemize}
\item \textbf{¿Existe un catálogo oficial de cursos de capacitación para servidores públicos del municipio?} Si es así, ¿quién nos lo puede proporcionar y en qué formato está?

\item \textbf{¿Cuáles son las categorías de cursos que se manejan?} Actualmente el sistema contempla: obligatorio, optativo y especializado. ¿Son correctas o se necesitan otras?

\item \textbf{¿Cuántos niveles de progresión debe tener un servidor?} Actualmente son 4 niveles. ¿Qué cursos corresponden a cada nivel? ¿Cuántos cursos debe completar un servidor para subir de nivel?

\item \textbf{¿Hay cursos que apliquen solo para cierto nivel de gobierno} (federal, estatal, municipal)? ¿O todos los cursos están disponibles para todos?
\end{itemize}

\subsection{2. Instituciones capacitadoras}

\begin{itemize}
\item \textbf{¿Cuáles son las instituciones que actualmente imparten capacitación a los servidores públicos del municipio?} Necesitamos: nombre, dirección, persona de contacto, teléfono y correo electrónico.

\item \textbf{¿Se trabaja con universidades, centros de capacitación municipales, plataformas en línea, o una combinación?}

\item \textbf{¿Las instituciones cambian frecuentemente o es un catálogo relativamente estable?}
\end{itemize}

\subsection{3. Servidores públicos existentes}

\begin{itemize}
\item \textbf{¿Existe un padrón actual de servidores públicos del municipio?} ¿En qué formato está (Excel, sistema anterior, nómina)?

\item \textbf{¿Qué datos del padrón ya están validados (RFC, CURP)?} El sistema valida ambos con formato oficial.

\item \textbf{¿Cuántos servidores públicos aproximadamente se van a registrar?} Esto nos ayuda a dimensionar la base de datos y los tiempos de carga.

\item \textbf{¿Las dependencias del municipio están estandarizadas en algún catálogo?} Nos serviría para un dropdown en el formulario en lugar de texto libre.
\end{itemize}

\subsection{4. Datos de contacto y responsables}

\begin{itemize}
\item \textbf{¿Quién será el administrador principal del sistema?} Nombre y correo electrónico para crear la cuenta inicial.

\item \textbf{¿Cuántos capturistas se necesitan y de qué áreas?} Para crear sus cuentas con el rol correspondiente.

\item \textbf{¿Hay personal de consulta (solo lectura) que necesite acceso?} Para asignarles el rol de consultor.
\end{itemize}

\newpage

\section{II. Preguntas Operativas --- Reglas de Negocio}

\subsection{5. Roles y permisos}

El sistema actualmente maneja 4 roles:

\vspace{8pt}

\begin{center}
\begin{tabular}{>{\bfseries}l p{10cm}}
\toprule
\textbf{Rol} & \textbf{Permisos} \\
\midrule
Administrador & Control total: usuarios, servidores, cursos, instituciones, solicitudes, reportes, auditoría \\
Capturista & Crear servidores, importar CSV, subir archivos \\
Consultor & Solo lectura: ver servidores, reportes, exportar datos \\
Usuario & Portal personal: ver cursos, solicitar inscripción, ver su progreso \\
\bottomrule
\end{tabular}
\end{center}

\vspace{8pt}

\begin{itemize}
\item \textbf{¿Estos roles son suficientes o se necesitan otros?} Ejemplo: supervisor de área, director, coordinador de capacitación.

\item \textbf{¿El capturista debería poder aprobar solicitudes de inscripción a cursos, o eso es exclusivo del administrador?}

\item \textbf{¿Los consultores pueden ver datos de todas las dependencias o solo de la suya?}
\end{itemize}

\subsection{6. Proceso de baja de servidores}

El sistema permite dos formas de baja:

\begin{enumerate}[leftmargin=2em]
\item \textbf{Baja directa por admin:} El administrador desactiva al servidor desde la gestión.
\item \textbf{Solicitud de baja por usuario:} El servidor solicita su propia baja con motivo, y el admin aprueba.
\end{enumerate}

\begin{itemize}
\item \textbf{Cuando un servidor fallece, ¿quién y cómo registra la baja?} ¿El admin directamente, o se requiere un documento soporte?

\item \textbf{¿Se necesitan motivos predefinidos para la baja?} Ejemplo: renuncia, jubilación, fallecimiento, cambio de adscripción, término de contrato. ¿O se deja como texto libre?

\item \textbf{¿Los datos de un servidor dado de baja se conservan indefinidamente o hay una política de retención/eliminación?}
\end{itemize}

\subsection{7. Niveles y clasificación de servidores}

Existe la referencia de que los servidores públicos se clasifican por niveles con valores numéricos (por ejemplo, nivel 5 y nivel 7), y que la asignación de cursos depende de dicho nivel. Necesitamos aclarar:

\begin{itemize}
\item \textbf{¿Existe una clasificación por niveles tabulares o jerárquicos de los servidores públicos?} (ejemplo: nivel 1 al 10, rangos salariales, categorías). ¿Cuántos niveles son y cuáles son?

\item \textbf{¿Los cursos de capacitación se asignan según este nivel?} Es decir, ¿un servidor de nivel 5 tiene acceso a ciertos cursos y uno de nivel 7 a otros diferentes?

\item \textbf{¿Este nivel es fijo (por puesto/tabulador) o cambia conforme el servidor avanza en capacitación?}

\item \textbf{¿Se nos puede proporcionar la tabla de niveles con los cursos que corresponden a cada uno?} Esto es necesario para configurar correctamente el catálogo en el sistema.

\item \textbf{¿Este nivel es independiente del nivel de gobierno (federal, estatal, municipal)?} ¿O se combinan ambos criterios para determinar qué cursos aplican?
\end{itemize}

\subsection{8. Proceso de capacitación}

\begin{itemize}
\item \textbf{¿Cómo se valida que un servidor completó un curso?} ¿El admin marca como completado manualmente, o se requiere evidencia (constancia, calificación)?

\item \textbf{¿Hay un tiempo máximo para completar un curso una vez aprobada la solicitud?}

\item \textbf{¿Un servidor puede solicitar el mismo curso más de una vez?} (ejemplo: si fue rechazado, ¿puede volver a solicitarlo?)

\item \textbf{¿Hay un máximo de cursos simultáneos que un servidor puede tener aprobados?}
\end{itemize}

\vfill

\begin{center}
\color{grissuave}\small\itshape
Documento para reunión de trabajo.\\
Las respuestas permitirán alimentar el sistema y definir reglas operativas.
\end{center}

\end{document}
